\documentclass[12pt,a4paper]{article}

\usepackage[T1]{fontenc}
\usepackage[utf8]{inputenc}
\usepackage{lmodern}
\usepackage{amsmath,amssymb,bm}
\usepackage{graphicx}
\usepackage{array}
\usepackage[a4paper,margin=2.5cm]{geometry}
\usepackage[hidelinks]{hyperref}
\usepackage{setspace}

\setstretch{1.15}
\setlength{\parindent}{0.7cm}
\setlength{\parskip}{0pt}
\setlength{\emergencystretch}{2em}
\newcolumntype{L}[1]{>{\raggedright\arraybackslash}p{#1}}

\title{Refractive Gravity: A Phase-Field Theory}
\author{}
\date{}

\begin{document}

\maketitle

\begin{abstract}
Refractive Gravity (RefG) is a phase-spatial effective field theory with an Einstein--Hilbert metric core and a medium sector built from one phase scalar $\Phi$ and three spatial scalars $\phi^A$. This paper establishes the gravitational core of the model: metric variation defines the active medium stress, the weak static Bianchi/Tolman--Oppenheimer--Volkoff balance gives the source-index chain, and minimal metric coupling recovers Newton's law for slow test bodies in the weak spherical limit. The weak-Solar branch is compatible with $\gamma=\beta=1$ and $q_{2{\rm PN}}=7/4$; the displayed preferred-frame source combination vanishes, and the principal tensor-speed sector has $c_g=c$.

The compact phase-spherical branch is treated as a separate static branch of the same normalized-invariant EFT sector. On this branch the phase-normalized $F_{\rm min}$ sector is stress-quiet, while the projected phase-pressure deficit source closes the field-equation residual of the exponential exterior. The resulting static branch has photon sphere $r_{\rm ph}=r_s$, critical impact parameter $b_c=e\,r_s$, innermost stable circular orbit radius $r_{\rm ISCO}=\varphi^2r_s$, and orbital frequency $0.931$ times the Schwarzschild value at the same mass. These are conditional static branch markers. Direct black-hole-imaging comparison requires the rotating exterior, finite-core dynamics, QNM/echo stability, and ray-traced plasma/accretion modeling.
\end{abstract}

\section{Introduction}

General Relativity (GR) supplies the metric core of gravitational dynamics. RefG keeps this core and adds a phase-spatial medium sector whose metric variation produces an active stress tensor. In the weak static limit that stress is read as a refractive profile.

The name ``Refractive Gravity'' refers to this source mechanism: the active stress of a phase-pressure deficit is encoded as a refractive metric response for matter and light. Matter remains minimally coupled to the metric; the effective-medium difference is that the stress-to-index chain recovers the Newtonian profile near a localized weak source, while the compact sector gives conditional static branch markers after a separate source/load selection.

The construction is an effective field theory written with the usual metric variables and scalar fields for the medium state. The phase field $\Phi$ tracks the temporal phase-pressure degree of freedom, while the three fields $\phi^A$ track the spatial body of the medium. Their invariants define the local energy functional. This follows the same conservative logic used in scalar-tensor and gravitational EFT work \cite{Horndeski1974,Gubitosi2013,BelliniSawicki2014,Deffayet2009,Kobayashi2011,Cheung2008}, with a different weak-field source: the refractive index produced by the active stress.

The observed equality of gravitational-wave and light speeds restricts the tensor sector \cite{Abbott2017GW,Abbott2017GRB,Creminelli2017,Sakstein2017}. RefG uses the Horndeski-compatible constant-$G_4$ sector, so the principal tensor speed is luminal. The spatial-medium part is related in spirit to solid and framid EFTs \cite{Endlich2013,Nicolis2015,Endlich2014,Bordin2017}; in RefG it serves as a gravitational source channel.

The compact static branch yields an exponential exterior. Similar exponential geometries have an independent literature, including analyses as traversable-wormhole geometries \cite{Boonserm2018}; strong-field alternatives to GR provide the comparison setting \cite{Damour1992,Damour1993,Damour1996}. In RefG this exterior is derived as the phase-spherical compact branch from the phase equation, biconformal map, and projected phase-pressure deficit source; the diagonal field-equation residual is checked explicitly below.

This paper assembles the checkable gravitational-core chain: action, active stress, refractive source, Newtonian limit, local principal-symbol stability, weak-Solar compatibility, tensor speed, inertial identity, and conditional compact static markers.

\medskip

\noindent\textbf{Terminology}\par\smallskip
\begin{list}{}{\setlength{\leftmargin}{0pt}\setlength{\itemindent}{0pt}\setlength{\labelwidth}{0pt}\setlength{\labelsep}{0pt}\setlength{\itemsep}{0pt}\setlength{\parsep}{0pt}\setlength{\parskip}{0pt}\setlength{\topsep}{0pt}\setlength{\partopsep}{0pt}}
\item \textbf{Base medium.} The phase-spatial EFT substrate represented by $\Phi$ and $\phi^A$; its invariants define the medium functional $F$.
\item \textbf{Active stress.} The source tensor obtained by metric variation of the medium action, $\Theta_{\mu\nu}^{\rm RefG}$, entering $G_{\mu\nu}=8\pi G\Theta_{\mu\nu}^{\rm RefG}$ in exterior regions.
\item \textbf{Refractive source.} The active pressure/deficit channel that fixes $h_{\rm eff}$ and $n_{\rm eff}$ through the Bianchi/TOV balance.
\item \textbf{Source-index chain.} The local calculation $\Pi_\Delta'\to h_{\rm eff}\to n_{\rm eff}\to a_r$.
\item \textbf{Phase-pressure deficit.} The projected quadratic channel $\Delta_P=Z_\Delta^\perp/(8\pi G)$ whose active stress supports the compact exterior.
\item \textbf{Phase-normalized strain invariants.} The $H$-rescaled variables $\widehat Y$, $\widehat\lambda_r$, and $\widehat\lambda_t$ used to evaluate $F_{\rm min}^{(H)}$ in a branch frame.
\item \textbf{Biconformal map.} The compact-branch map $h=-\phi/2$ with paired metric factors $e^{-2h}$ and $e^{2h}$.
\item \textbf{Branch.} A background regime of the same normalized-invariant EFT sector, selected by compactness, boundary matching, exterior source weight, and source amplitude.
\item \textbf{Weak-Solar branch.} The unloaded weak-field $H=0$ branch with $\gamma=\beta=1$, vanishing preferred-frame source, and $q_{2{\rm PN}}=7/4$.
\item \textbf{Compact phase-spherical branch.} The loaded static $H=H_\Delta=h$ branch in which the reduced phase equation and projected deficit source give the exponential exterior.
\item \textbf{Dynamic completion.} The local principal-symbol, static-silent $\epsilon_B,\epsilon_M$ extension that changes the scalar-longitudinal kinetic/mixing entries and opens the subluminal window; full curved-background nonlinear hyperbolicity remains a separate gate.
\item \textbf{ESS operator.} The static-silent effective-solid-sector operator $\delta_{AB}(u^\mu\partial_\mu\phi^A)(u^\nu\partial_\nu\phi^B)$ used to give the Solar solid mode positive kinetic energy.
\item \textbf{Raw-invariant diagnostic.} The compact-exterior check using unnormalized raw invariants; it identifies a diagnostic obstruction, not the compact action.
\end{list}

\subsection*{Validation Status and Scope}

The claims below are used only in the scope in which the algebraic checks close.
A closed entry means that the stated calculation is closed in that sector; the
right column lists what is not concluded from it.
\begin{center}
\small
\renewcommand{\arraystretch}{1.18}
\begin{tabular}{L{0.22\textwidth}L{0.35\textwidth}L{0.33\textwidth}}
\hline
Sector & Article-level status & Scope boundary \\
\hline
Refractive source and Newton limit &
Action stress gives the weak static source-index chain and Newton's law in the
weak spherical limit. &
Full nonlinear, non-spherical medium PDEs and the microscopic origin of $G$ are
not claimed here. \\
Weak-Solar branch &
The $1$PN branch has $\gamma=\beta=1$ and $\alpha_i=0$ in the minimal chain;
$q_{2{\rm PN}}=7/4$ is the weak-Solar branch value. &
The full raw Solar-System likelihood and the complete PPN table remain outside
this paper. \\
Local stability &
The scalar-longitudinal sector has a healthy local principal symbol after the
static-silent dynamic repair; a representative subluminal point is displayed. &
Full curved-background nonlinear hyperbolicity and constraint closure remain
open. \\
Tensor speed &
The local tensor sector has constant $G_4$, $G_5=0$, and $c_g=c$. &
Full cosmological wave propagation, damping, and waveform/catalog tests remain
open. \\
Compact phase branch &
The static exponential exterior source ledger closes: phase-normalized
$F_{\rm min}$ is stress-quiet, the projected deficit closes the diagonal
residual, and the source amplitude selects a static finite-core window. &
No black-hole replacement, EHT support, rotating exterior, QNM/echo stability,
or plasma/accretion ray-tracing claim is made here. \\
\hline
\end{tabular}
\end{center}

\section{The Base Medium and the Physical Picture}

RefG starts from a base medium that defines operational distance, time, mass, and wave propagation as local regimes of one substrate. A homogeneous state is self-calibrated by rods, clocks, and atoms formed within the same local phase-pressure state. Observable gravity begins when the medium state becomes spatially nonuniform.

One base medium supports distinct mechanical responses. The phase, pressure, and spatial deformation channels used in this paper are separate responses of the same substrate, linked by internal medium interactions. A localized energy concentration creates a deficit in the phase pressure of the base medium. In the weak static reading, geodesic motion can be rewritten as motion in the gradient of the resulting refractive index, giving Newton's force in the weak spherical regime. The compact phase-dominated exterior is selected by a separate source/load condition and gives the exponential static branch.

\section{Fields, Invariants, and the Action}

\subsection{Invariant Basis and Action}

The fundamental gravitational variables are the metric $g_{\mu\nu}$, the phase scalar $\Phi$, and three spatial medium scalars $\phi^A$, $A=1,2,3$. The metric signature is $(+---)$. The medium invariants are
\begin{align*}
Y &= g^{\mu\nu}\partial_\mu\Phi\,\partial_\nu\Phi,\\
B^{AB} &= -g^{\mu\nu}\partial_\mu\phi^A\,\partial_\nu\phi^B,\\
I_1 &= {\rm Tr}\,B,\qquad
I_2=\frac{1}{2}\left[({\rm Tr}\,B)^2-{\rm Tr}(B^2)\right],\qquad
I_3=\det B.
\end{align*}
In this normalization $\Phi$ is a phase-clock scalar and $\phi^A$ are spatial medium coordinate scalars. On the unit solid background one has $\phi^A=x^A$ and $Y=1$, $B^{AB}=\delta^{AB}$. In the convention $\mathcal L=M_*^4F$, the invariants are dimensionless and the coefficients $c_i$ are effective medium response parameters; the physical density dimension is carried by $M_*^4$.

The fields $\Phi$ and $\phi^A$ represent different operational channels of the same base medium. The $YI_1$ term couples these channels while keeping the phase response and the spatial response as separate dynamical properties. This distinction is central in the scalar-longitudinal stability criterion below.

The minimal medium core used in this paper is
\begin{align*}
F_{\rm min}={}&
c_Y Y+c_{Y2}Y^2+c_{I1}I_1+c_{I1sq}I_1^2+c_{I2}I_2+c_{I3}I_3
+c_{YI1}YI_1.
\end{align*}
The action uses this polynomial as a local-frame functional.  Let
$F_{\rm min}^{(H)}$ denote the same polynomial evaluated on the
$H$-normalized strain invariants defined below.  On the weak-Solar frame
$H=0$, $F_{\rm min}^{(H)}$ reduces to the raw invariant expression displayed
above.
The article-level leading action is
\begin{align*}
S=\int d^4x\sqrt{-g}\left[
\frac{M_{\rm Pl}^2}{2}R+M_*^4 F_{\rm min}^{(H)}
+\mathcal L_{\rm EFT}^{>{\rm min}}
\right]+S_{\rm m}[g_{\mu\nu},\psi].
\end{align*}
The matter sector is minimally coupled to the metric. The term
$\mathcal L_{\rm EFT}^{>{\rm min}}$ is the higher EFT sector of the same
medium. The compact-exterior projected phase-pressure deficit operator is
written with an independent deficit scalar $H_\Delta$:
\[
u_\mu=\frac{\partial_\mu\Phi}{\sqrt Y},\qquad
\gamma^{\mu\nu}=u^\mu u^\nu-g^{\mu\nu},
\]
\[
Z_\Delta^\perp=
\gamma^{\mu\nu}\partial_\mu H_\Delta\,\partial_\nu H_\Delta,
\]
\[
\mathcal L_{\rm EFT}^{>{\rm min}}\supset
\mathcal L_\Delta^\perp,\qquad
\mathcal L_\Delta^\perp=
\omega_\Delta\,\frac{Z_\Delta^\perp}{8\pi G}.
\]
Here $\omega_\Delta$ is a branch-selected dimensionless EFT coefficient; the
compact branch normalization used below sets $\omega_\Delta=1$. In the
weak-Solar branch this channel is treated as unloaded at the leading order used
here. In the compact $C^2$ matching branch
$H_\Delta$ carries the matched exterior source law through
$\mathcal L_\Delta^\perp$, while $I_3$ remains an internal medium invariant.
The compact calculation below is this leading EFT truncation:
$\mathcal L_\Delta^\perp$ is the first active projected-deficit operator on the
matched pure-phase exterior, while other higher-sector operators are left as
beyond-core corrections unless a later branch derivation makes them active.

\subsection{Branch Frames and Phase Normalization}

The compact branch uses the same $F_{\rm min}^{(H)}$ polynomial in a different
local frame.  Let $H$ be the phase-normalization field and define the
phase-normalized
strain invariants
\[
\widehat Y=e^{-2H}Y,\qquad
\widehat\lambda_r=e^{2H}\lambda_r,\qquad
\widehat\lambda_t=e^{2H}\lambda_t.
\]
The hatted invariants $\widehat I_n$ are formed from
$\widehat\lambda_r$ and $\widehat\lambda_t$ in the same way that $I_n$ are
formed from $\lambda_r$ and $\lambda_t$.
Notation is fixed as follows: $h_{\rm eff}$ is the weak static source profile,
$h$ is the compact metric phase, $H$ normalizes the $F_{\rm min}$ strain
invariants, and $H_\Delta$ is the independent deficit-source scalar. On the
compact pure-phase exterior $H=H_\Delta=h$; in the weak-Solar branch condition below
$H=0$.
These equalities are branch and matching conditions, not universal field
identities: $H$ is the frame variable used to evaluate $F_{\rm min}$, while
$H_\Delta$ is the independent higher-sector source field, and away from the
compact pure-phase exterior they need not coincide.
In the displayed core $H$ is auxiliary: no derivative of $H$ occurs in
$F_{\rm min}^{(H)}$, and the compact-source gradients are carried by
$H_\Delta$ through $\mathcal L_\Delta^\perp$. Variation with respect to $H$
therefore gives the algebraic branch constraint
\[
S_H\equiv\frac{\partial F_{\rm min}^{(H)}}{\partial H}=0.
\]
The weak-Solar diagnostic below verifies this constraint for $H=0$ through the
displayed Solar order, while the compact pure-phase branch satisfies it exactly
by the tadpole-free unit state. Thus the two exteriors are self-consistent
solutions of the same auxiliary-$H$ constraint; the finite-core load/source
rule below fixes the selected compact exterior window.
Branches are background regimes of the same normalized-invariant EFT sector:
the weak-Solar branch has the unloaded normalization $H=0$, while the compact
pure-phase exterior has the loaded normalization $H=h$. The raw
absolute-invariant compact insertion checked below is a diagnostic substitution, not a
replacement of the action used for the compact branch.
In the compact pure-phase exterior,
\[
H=h,\qquad Y=e^{2h},\qquad
\lambda_r=\lambda_t=e^{-2h},
\]
and therefore
\[
\widehat Y=\widehat\lambda_r=\widehat\lambda_t=1.
\]
The Solar-family $F_{\rm min}$ slice used below is tadpole-free at this unit
state:
\[
F_{\rm min}(1,1,1)=0,\qquad
\left.\frac{\partial F_{\rm min}}{\partial\widehat Y}\right|_1=
\left.\frac{\partial F_{\rm min}}{\partial\widehat\lambda_r}\right|_1=
\left.\frac{\partial F_{\rm min}}{\partial\widehat\lambda_t}\right|_1=0.
\]
Consequently the compact phase-normalized $F_{\rm min}$ sector has
\[
\widehat\Theta_F^t{}_t=
\widehat\Theta_F^r{}_r=
\widehat\Theta_F^\theta{}_\theta=
\widehat\Theta_F^\phi{}_\phi=0,\qquad
E_f[\widehat F_{\rm min}]=0,
\]
on the pure-phase exterior. This is an action-level statement: the compact
branch varies the phase-normalized $F_{\rm min}$ action from the start. The raw
absolute-invariant compact insertion gives the nonzero tensor checked below. If
the compact phase field is unloaded, $H=0$, the hatted invariants reduce to the
raw weak-branch invariants.

\subsection{Weak-Solar Diagnostic}

The weak-Solar branch condition is explicit. For $U=r_s/r$,
\[
Y=1+U+O(U^2),\qquad
\lambda_r=1-U+O(U^2),\qquad
\lambda_t=1+O(U^2).
\]
Raw $F_{\rm min}$ and phase-normalized $F_{\rm min}$ with $H=0$ both give
\[
\Theta_F^t{}_t=\Theta_F^r{}_r=\Theta_F^\theta{}_\theta=0
\quad {\rm at}\quad O(U),
\]
and the augmented Solar ansatz $F_{\rm map}/r=1+\sigma U^2$ keeps the
exact-GR branch $\sigma=-1/2$. The determinant-locked trial
$H=-\ln(\lambda_r\lambda_t^2)/6$ gives instead
\[
\Theta_F^t{}_t=M_*^4c_{Y2}\frac{8}{3}U,\qquad
\Theta_F^r{}_r=\Theta_F^\theta{}_\theta=M_*^4c_{Y2}\frac{8}{9}U,
\]
and has no constant-strain exact-GR $2$PN solution. Thus $H$ is the
pressure/energy-deficit channel; the compact source is the projected deficit
operator.

The $H$ Euler source also vanishes on the weak-Solar branch. Write
\[
Y=1+y_1U+y_2U^2,\qquad
\lambda_r=1+r_1U+r_2U^2,\qquad
\lambda_t=1+t_1U+t_2U^2.
\]
For the same action template, with
$S_H\equiv\partial F_{\rm branch}/\partial H$ evaluated at $H=0$,
\[
\frac{S_H}{c_{Y2}}
=8(y_1+r_1+2t_1)U
+4\left(r_1^2+20r_1t_1+2r_2+12t_1^2+4t_2-y_1^2+2y_2\right)U^2
+O(U^3).
\]
The weak-Solar strain has $y_1=1$, $r_1=-1$, $t_1=0$. On the exact-GR $2$PN
constant-strain branch one has $y_2=1$, $r_2=1$, $t_2=-1$. Hence
$S_H=0+O(U^3)$. Thus the unloaded branch $H=0$ is an $H$-equation solution
through the Solar order displayed here, rather than merely a formal reduction of the action.
The cancellation is structural: the first-order source is proportional to the
weak variation of $Y+I_1$, and the second-order source vanishes under the same
exact-GR vacuum condition that removes the $F_{\rm min}$ gradient and the
exterior $2$PN metric stress.
On the compact pure-phase branch the hatted invariants equal one, so the same
$H$ source vanishes exactly.

At weak static order the calculation chain from the minimal core to the source profile is direct. Metric variation of the action gives the active stress $\Theta^{\rm RefG}_{\mu\nu}$; the radial balance of this stress enters the Bianchi/Tolman--Oppenheimer--Volkoff (TOV) identity; the resulting pressure source determines the profile $h_{\rm eff}$; the slow-body source index reads that profile as $n_{\rm eff}=\exp(h_{\rm eff})$. Thus the weak source-index chain is the active medium response of the action stress, read through the differential TOV/Bianchi source equation.

The same action contains the technical background channels used for local stability diagnostics and for the weak-Solar branch. The combination
\[
C_6=Y+2I_1-I_3,
\]
and the invariant
\[
Z=\delta_{AB}(u^\mu\partial_\mu\phi^A)(u^\nu\partial_\nu\phi^B),
\qquad
u_\mu=\frac{\partial_\mu\Phi}{\sqrt{Y}},
\]
enter as a local completion of the medium background. The isolated $C_6/Z$
block is a diagnostic luminal completion; the scalar-longitudinal sector
exported below uses the static-silent dynamic repair at local principal-symbol
level.
The phase and spatial responses remain independent channels.

The active RefG stress is defined by
\[
\Theta^{\rm RefG}_{\mu\nu}
=-2\frac{\partial\mathcal L_{\rm RefG}}{\partial g^{\mu\nu}}
+g_{\mu\nu}\mathcal L_{\rm RefG}.
\]
Equivalently,
\[
\delta\!\int d^4x\sqrt{-g}\,\mathcal L_{\rm RefG}
=-\frac12\int d^4x\sqrt{-g}\,
\Theta^{\rm RefG}_{\mu\nu}\delta g^{\mu\nu}.
\]
With this convention the exterior field equation is
\[
G_{\mu\nu}=8\pi G\,\Theta^{\rm RefG}_{\mu\nu}
\qquad (T_{\mu\nu}^{\rm m}=0).
\]
This active stress enters the source-index chain below. The ordinary matter tensor $T_{\mu\nu}^{\rm m}$ and the medium source $\Theta_{\mu\nu}^{\rm RefG}$ are distinct objects. In an exterior region $T_{\mu\nu}^{\rm m}=0$, and the field-equation residual is read with the active medium source. A negative mixed time component of $\Theta^\mu{}_\nu$ in the compact branch denotes a phase-pressure deficit of the medium.

In the pure phase projection the model is in the minimal Horndeski-compatible class. With $Y=-2X$,
\[
G_2=-2c_YX+4c_{Y2}X^2,\qquad
G_3=0,\qquad
G_4=\frac{M_{\rm Pl}^2}{2},\qquad
G_5=0.
\]
Because $G_4$ is constant and $G_5=0$, the local principal tensor speed is
luminal in this tensor sector. This is not used as a full cosmological
wave-propagation calculation.

\section{The Refractive Source and Newton's Law}

The weak static metric is written as
\[
g_{tt}=\exp(-2h_{\rm eff}).
\]
For slow test matter the corresponding source index is
\[
n_{\rm eff}=\exp(h_{\rm eff}).
\]
The effective potential per unit mass is
\[
\frac{V_{\rm eff}}{m}=-c^2h_{\rm eff},
\]
and the radial acceleration is
\[
a_r=c^2\frac{d\ln n_{\rm eff}}{dr}
=c^2h_{\rm eff}'.
\]
The attraction is inward when $h_{\rm eff}$ decreases outward.

The source of $h_{\rm eff}$ is the active medium stress. For a static spherical configuration define the anisotropic pressure contrast
\[
\Delta_p=p_{\rm tan}-p_{\rm rad}.
\]
The Bianchi/TOV channel gives
\[
p_{\rm rad}'+\rho_{\rm eff}\Phi_N'
-\frac{2\Delta_p}{r}=0,
\qquad
\Phi_N=-c^2h_{\rm eff}.
\]
Equivalently,
\[
h_{\rm eff}'=
\frac{p_{\rm rad}'-2\Delta_p/r}{c^2\rho_{\rm eff}}.
\]
This is the linear source equation. It locates the active pressure gradient in the refractive profile.

The deficit part of the same source is quadratic in the refractive gradient:
\[
\Pi_\Delta=\frac{e^{-2h_{\rm eff}}(h_{\rm eff}')^2}{8\pi G}.
\]
Its associated effective density is
\[
\rho_\Delta=\frac{\Pi_\Delta'}{c^2h_{\rm eff}'}
=\frac{e^{-2h_{\rm eff}}\left(h_{\rm eff}''-(h_{\rm eff}')^2\right)}
{4\pi G c^2}.
\]
The working object is the RefG quadratic deficit channel:
the source contains the square of the refractive-profile gradient, and the
radial balance uses its derivative.
Thus the calculation chain is
\[
\Pi_\Delta' \longrightarrow h_{\rm eff} \longrightarrow n_{\rm eff}
\longrightarrow a_r.
\]

In the local spherical weak branch the anisotropic term is absent at leading order,
\[
\Delta_p=0,
\]
and the exterior source-free equation is
\[
\left(r^2h_{\rm eff}'\right)'=0.
\]
The Gauss normalization of the localized source charge fixes
\[
-r^2h_{\rm eff}'=\frac{GM}{c^2}=\frac{r_s}{2},
\]
and therefore
\[
h_{\rm eff}(r)=\frac{GM}{c^2r}=\frac{r_s}{2r},
\qquad r_s=\frac{2GM}{c^2}.
\]
Therefore
\[
\Phi_N=-c^2h_{\rm eff}=-\frac{GM}{r},
\qquad
a_r=c^2h_{\rm eff}'=-\frac{GM}{r^2}.
\]
Newton's law follows as the weak spherical limit of the refractive source-index chain.

\section{Local Stability}

\subsection{Bare Principal Symbol}

The local perturbation calculation is performed around the homogeneous medium background. The relevant coefficient combination is
\[
\Sigma=c_Y+3c_{YI1}.
\]
The phase and longitudinal scalar kinetic entries are
\[
K_{\Phi\Phi}=\Sigma+6c_{Y2},
\qquad
K_{\pi\pi}=-\Sigma-2c_{Y2}.
\]
The nondegenerate bare-minimal no-ghost window is
\[
c_{Y2}>0,\qquad
-6c_{Y2}<\Sigma<-2c_{Y2}.
\]

The coupled scalar principal symbol has determinant
\[
\det M(s)=p_2s^2+p_1s+p_0,\qquad s=\omega^2/k^2.
\]
Hyperbolicity and positive propagation speeds require
\[
p_2>0,\qquad p_1<0,\qquad p_0>0,\qquad
p_1^2-4p_2p_0\ge 0.
\]
The Solar-compatible family is given by
\begin{align*}
c_{I1}&=4c_{Y2}+2c_{YI1},\\
c_{I1sq}&=c_{Y2},\\
c_{I2}&=-10c_{Y2}-3c_{YI1},\\
c_{I3}&=8c_{Y2}+4c_{YI1},\\
c_Y&=-4c_{Y2}-2c_{YI1}.
\end{align*}
The physical Solar $2$PN branch imposes
\[
c_{YI1}=2c_{Y2}.
\]
On this slice
\[
c_Y=-8c_{Y2},\qquad
K_{\Phi\Phi}=4c_{Y2}>0,\qquad
K_{\pi\pi}=0.
\]
Thus the physical Solar point lies on the longitudinal boundary of the
bare-minimal window: the phase kinetic term is healthy, while the longitudinal
solid kinetic term is degenerate and must be read in the completed principal
symbol below.

\subsection{Static Completion Diagnostic}

The technical $C_6/Z$ background completion is first used as a diagnostic that
closes the flat kinetic matrix without changing the tensor sector. With
\[
\Phi=t+\chi,\qquad \phi^A=x^A+\pi^A,
\]
one has
\[
\delta C_6=2(\dot\chi+\partial_i\pi^i),
\qquad
Z=(\dot\pi_i-\partial_i\chi)^2.
\]
The quadratic completion is
\[
\mathcal L_2=\lambda_6(\delta C_6)^2+c_Z Z.
\]
For the Solar branch one first combines the $F_{\rm min}$ block with the
$C_6/Z$ block.  The scalar-longitudinal determinant is then read from the
following Fourier coefficients:
\[
A=4(c_{Y2}+\lambda_6),\quad B=c_Z,\quad C=c_Z,\quad
D=4(c_{Y2}+\lambda_6),\quad
M=8c_{Y2}-4\lambda_6-c_Z.
\]
Therefore
\[
\det M_\odot(s)=p_2s^2+p_1s+p_0,
\]
with
\begin{align*}
p_2&=4c_Z(c_{Y2}+\lambda_6),\\
p_1&=96c_{Y2}\lambda_6-8c_Z\lambda_6-48c_{Y2}^2+16c_{Y2}c_Z,\\
p_0&=4c_Z(c_{Y2}+\lambda_6).
\end{align*}
One explicit Solar kinetic point is
\[
c_{Y2}=1,\qquad \lambda_6=\frac14,\qquad c_Z=1.
\]
At this point
\[
\det M_\odot(s)=5(s-1)^2,
\]
so at this boundary point the two scalar-longitudinal roots are real and
luminal:
\[
s_1=s_2=1.
\]
This calculation fixes the status of the pre-completion $F_{\rm min}+C_6/Z$ sector.
Since $p_0=p_2$, the product of the two roots is always one.  Hence two
distinct real roots cannot both be subluminal: if one is below one, the other
is above one.  The representative point $5(s-1)^2$ is a boundary luminal
double root. The completed operator below supplies the open subluminal window.
In this pre-completion form the $C_6/Z$ block overconstrains the phase lag and the
longitudinal medium response.

\subsection{Dynamic Completion}

The dynamic completion separates the phase-spatial lag into an independent
channel.  The original $Z$ invariant measures the motion of the medium
relative to the phase normal,
\[
U^A=u^\mu\partial_\mu\phi^A.
\]
In the local scalar-longitudinal limit this reads
\[
U_L=\dot\pi_L-\nabla_L\chi.
\]
It therefore contains the longitudinal medium velocity and the phase tilt in
one combination.  The completed channel separates them by defining
\[
Q^A=-g^{\mu\nu}\partial_\mu\phi^A\,\partial_\nu\Phi,\qquad
W^A=U^A+Q^A,
\]
so locally $Q_L=\nabla_L\chi$ and $W_L=\dot\pi_L$.  The following operator is
static-silent: it vanishes for a static source but contributes to the
principal symbol,
\[
\Delta\mathcal L_{\rm dyn}
=\epsilon_B W_A W^A+2\epsilon_M W_A Q^A.
\]
For a static source $\dot\pi_L=0$, hence $W^A=0$, so this operator leaves the
static Solar gradient core unchanged.  The principal-symbol dynamic entries
are shifted as
\[
A_{\rm new}=4(c_{Y2}+\lambda_6),\qquad
B_{\rm new}=c_Z+\epsilon_B,\qquad
C_{\rm new}=c_Z,
\]
\[
D_{\rm new}=4(c_{Y2}+\lambda_6),\qquad
M_{\rm new}=8c_{Y2}-4\lambda_6-c_Z+\epsilon_M.
\]
The corresponding $(\chi,\pi_L)$ principal matrix is
\[
\mathcal P(\omega,k)=
\begin{pmatrix}
A_{\rm new}\omega^2+C_{\rm new}k^2 & M_{\rm new}\omega k\\
M_{\rm new}\omega k & B_{\rm new}\omega^2+D_{\rm new}k^2
\end{pmatrix}.
\]
This matrix follows directly from the quadratic action.  The original $Z$ term gives
$(\omega\pi_L-k\chi)^2$, while the dynamic operator gives
\[
\Delta\mathcal L_{\rm dyn}
=\epsilon_B\omega^2\pi_L^2+2\epsilon_M\omega k\chi\pi_L.
\]
Thus $\epsilon_B$ enters only the longitudinal kinetic entry $B_{\rm new}$,
and $\epsilon_M$ enters only the mixing entry $M_{\rm new}$.  The entries
$A_{\rm new}$, $C_{\rm new}$, and $D_{\rm new}$ remain the static-core entries.

Here $\epsilon_B$ is the independent longitudinal-medium dynamic response and
$\epsilon_M$ is the independent phase-spatial lag contribution.  Since $C_{\rm
new}$ and $D_{\rm new}$ are unchanged, the static Solar gradient core is
unchanged.  Only the dynamic condition that produced the pre-completion $p_0=p_2$ locking is
changed.
The completed EFT treats $\lambda_6$, $c_Z$, $\epsilon_B$, and $\epsilon_M$ as
medium response coefficients and admits an open healthy flat principal-symbol
window; the numerical point below is representative.

The entries $C_{\rm new}$ and $D_{\rm new}$ already contain the full flat
gradient sector on the physical Solar slice. The $Y$, $I_1$, $I_2$, $I_3$,
and $YI_1$ terms are first inserted in the full $F_{\rm min}$ principal
symbol. On the physical Solar slice the full $F_{\rm min}$ principal-symbol
entries are
\[
A_F=4c_{Y2},\qquad B_F=0,\qquad C_F=0,\qquad
D_F=4c_{Y2},\qquad M_F=8c_{Y2}.
\]
The $C_6/Z$ completion then adds $C=c_Z$ and $D=4\lambda_6$, and
$\Delta\mathcal L_{\rm dyn}$ changes only $B$ and $M$.

The completed determinant is
\[
\det M_{\rm rep}(s)=(A_{\rm new}s+C_{\rm new})(B_{\rm new}s+D_{\rm new})
-M_{\rm new}^2s.
\]
At the representative point
\[
c_{Y2}=1,\qquad \lambda_6=\frac14,\qquad c_Z=1,\qquad
\epsilon_B=1,\qquad \epsilon_M=\sqrt{\frac{83}{2}}-6,
\]
one obtains
\[
\det M_{\rm rep}(s)=\frac{20s^2-29s+10}{2}.
\]
Its roots are
\[
s_1=\frac{29-\sqrt{41}}{40}\simeq0.5649,\qquad
s_2=\frac{29+\sqrt{41}}{40}\simeq0.8851.
\]
On the same Solar slice $c_{Y2}=1$, $\lambda_6=1/4$, $c_Z=1$, the completed
determinant is
\[
\det M_{\rm rep}(s)
=5(1+\epsilon_B)s^2+\left[26+\epsilon_B-(6+\epsilon_M)^2\right]s+5.
\]
With $M_2=(6+\epsilon_M)^2$, real, positive, strictly subluminal roots are
obtained in the open window
\[
26+\epsilon_B+10\sqrt{1+\epsilon_B}<M_2<36+6\epsilon_B.
\]
The width of this window is
\[
5\left(\sqrt{1+\epsilon_B}-1\right)^2,
\]
so it is finite for every $\epsilon_B>0$ and collapses only at
$\epsilon_B=0$.  At the representative value $\epsilon_B=1$ this gives
$41.142\ldots<M_2<42$, and the chosen value $M_2=41.5$ lies inside it.

Both roots are real, positive, and strictly subluminal.  Thus the Solar
$2$PN exterior is unchanged, the phase kinetic term is positive, the
longitudinal solid kinetic term is completed, and the scalar-longitudinal speed
condition is satisfied by an independent dynamic channel.  The pre-completion $F_{\rm min}+C_6/Z$
block is a diagnostic of the overconstraint; the principal symbol of
$\Delta\mathcal L_{\rm dyn}$ gives the subluminal window.
Because the two completed roots are distinct, the pre-completion luminal double
root is confined to the diagnostic sector. At the flat principal-symbol level
the completed sector is nondegenerate; curved-background hyperbolicity remains a
separate technical problem.

\begin{figure}[t]
\centering
\includegraphics[width=0.78\textwidth]{figures/refg_stability_window.pdf}
\caption{Completed scalar-longitudinal principal-symbol window.  The shaded
region is the locally healthy subluminal interval
$26+\epsilon_B+10\sqrt{1+\epsilon_B}<M_2<36+6\epsilon_B$, where
$M_2=(6+\epsilon_M)^2$.  The point
$(\epsilon_B,M_2)=(1,41.5)$ gives
$s_1=(29-\sqrt{41})/40\simeq0.565$ and
$s_2=(29+\sqrt{41})/40\simeq0.885$.}
\label{fig:stability-window}
\end{figure}

\subsection{Solar Branch Link}

This completion is the link to the Solar $2$PN branch. The Solar
$q_{2{\rm PN}}=7/4$ branch below uses
\[
c_{YI1}=2c_{Y2}.
\]
On the Solar $1$PN coefficient family one has
\[
\Sigma=-4c_{Y2}+c_{YI1},
\]
and therefore the bare minimal core gives
\[
K_{\pi\pi}=2c_{Y2}-c_{YI1}.
\]
On the Solar $2$PN condition this bare factor is zero. This is a degeneracy of
the minimal polynomial core. In the completed kinetic frame the $C_6/Z$
channel is replaced, for the dynamical completion, by the static-silent
effective-solid-sector (ESS) operator
\[
\delta_{AB}(u^\mu\partial_\mu\phi^A)(u^\nu\partial_\nu\phi^B),
\]
which vanishes on the background and in the first static source variation while
adding positive solid kinetic energy in the quadratic perturbation action. Thus
the Solar static $2$PN exterior is unchanged, and the local scalar-longitudinal
sector is governed by the completed determinant above.  The displayed completed
point gives two real subluminal roots.

\section{The Solar System}

The Solar-System branch must reproduce the standard post-Newtonian (PN) tests \cite{Will1993,Will2014,Bertotti2003}. The spherical calculation is first written in a Schwarzschild-like areal-radius working coordinate. With $u=GM/(c^2R)$, the areal time coefficient is denoted by $\beta_A$:
\begin{align*}
g_{tt}^{A}&=1-2u+2(\beta_A-1)u^2+\cdots,\\
g_{RR}^{A}&=-(1+2\gamma u+a_2u^2+\cdots).
\end{align*}
The standard isotropic parameterized post-Newtonian (PPN) time metric is
\[
g_{tt}^{\rm PPN}=1-2U+2\beta_{\rm PPN}U^2+\cdots.
\]
The areal and isotropic radii obey
\[
R=\rho(1+\gamma U+O(U^2)),\qquad
u=U-\gamma U^2+O(U^3),
\]
and therefore
\[
\beta_{\rm PPN}=\beta_A+\gamma-1.
\]
On the Solar first post-Newtonian branch $\gamma=1$, so $\beta_{\rm PPN}=\beta_A$. The first post-Newtonian matching gives
\[
\gamma=1.
\]
At second order, with
\[
a_2=4,
\]
the same branch gives
\[
\beta_A=1,
\qquad
\beta_{\rm PPN}=1.
\]
The same areal-to-isotropic matching gives $b_2=(a_2-\gamma^2)/2$; hence
$a_2=4$ and $\gamma=1$ imply $b_2=3/2$.
The compact exponential branch has the isotropic time factor
\[
\exp(-2U)=1-2U+2U^2+\cdots,
\]
so its standard PPN time coefficient is also $\beta_{\rm PPN}=1$. The Solar branch and the compact phase branch share the first PPN $\beta$ matching and separate at the 2PN/strong-field branch level.
In the displayed static scalar reduction, the preferred-frame source
combination is the coefficient
\[
K_{\rm pf}
=-c_{I1}-6c_{I1sq}-2c_{I2}-c_{I3}
+c_Y+2c_{Y2}+2c_{YI1}.
\]
On the Solar-compatible family,
\[
K_{\rm pf}=0.
\]
This is the local coefficient condition used here for
\[
\alpha_1=\alpha_2=\alpha_3=0,
\]
with the full vector-sector PN analysis belonging to the complete PPN extension
\cite{WillNordtvedt1972,Kostelecky2022}.

The second post-Newtonian optical coefficient used by the Solar branch is the
physical weak-Solar branch value:
\[
q_{2{\rm PN}}=\frac{7}{4}.
\]
Restricted zero-stress and dark-energy-scale checks test this branch separation;
they are not a full raw Solar-System likelihood. The compact phase branch is a
separate strong-field exterior; its conditional static markers are the shadow
and innermost stable circular orbit (ISCO) quantities derived below.

The Solar branch is obtained by imposing the second-order relation
\[
c_{YI1}=2c_{Y2}.
\]
The radial medium deformation $s(r)$ satisfies
\[
2rs'(r)+2s(r)+1=0,
\]
so that
\[
s(r)=\frac{C}{r}-\frac12.
\]
The term $C/r$ is a boundary tail; the local second-order coefficient is fixed by the constant part $-1/2$. With this substitution the angular residual becomes the radial-deformation balance, and the metric part is the GR-compatible isotropic $2$PN exterior.

On the same condition the bare $F_{\rm min}$ longitudinal kinetic prefactor is
\[
K_{\pi\pi}=2c_{Y2}-c_{YI1}=0.
\]
The physical Solar $2$PN branch is therefore paired with the completed kinetic
frame described in the stability section. The ESS kinetic completion vanishes
on the static background and in the static field equations, leaving unchanged
the equation for $s(r)$, the choice $a_2=4$, or the result
$q_{2{\rm PN}}=7/4$. It contributes only at perturbation level, where it gives
the solid mode positive kinetic energy and satisfies the no-ghost and speed criteria.

Write
\[
U=\frac{GM}{c^2\rho},
\qquad
B=1-2U+2\beta U^2+O(U^3),
\qquad
A=1+2\gamma U+b_2U^2+O(U^3).
\]
The static light index is
\[
n_{\rm light}=\sqrt{\frac{A}{B}},
\]
therefore
\[
n_{\rm light}=1+(\gamma+1)U+
\left(\frac{b_2}{2}-\beta-\frac{\gamma^2}{2}
+\gamma+\frac32\right)U^2+O(U^3).
\]
On the physical Solar branch
\[
\gamma=1,\qquad \beta=1,\qquad b_2=\frac32,
\]
and hence
\[
n_{\rm light}=1+2U+\frac74U^2+O(U^3).
\]
Thus $q_{2{\rm PN}}=7/4$. Cassini and lunar-laser-ranging tests constrain this
weak-Solar metric branch, but a full raw Solar-System likelihood is a separate
data-analysis step. RefG changes the source mechanism, while the displayed
weak-Solar $1$PN/$2$PN optical coefficients remain GR-compatible.

The technical background completion is negligible on Solar scales. Its dimensionless contribution is
\[
\Delta_{\rm Solar}\sim \Lambda R_\odot^2\sim 5.3\times10^{-35}.
\]
The post-Newtonian Solar matching is unchanged.

The projected deficit operator is also directly small in the weak-Solar
exterior. With compactness $C=r_s/R$, its local source scale relative to the
leading Newtonian curvature scale is
\[
\frac{D_\Delta}{D_N}=\frac{C e^{-C}}{4}.
\]
At the Solar surface this is about $1.06\times10^{-6}$. Thus the diffuse
medium stress dominates the weak-Solar boundary matching over the
phase-dominated deficit exterior. In the compact $C^2$ matching branch the same
projected operator, with the compact-branch normalization $\omega_\Delta=1$,
balances the exponential exterior.

\section{Tensor Waves}

The tensor sector follows from the Horndeski-compatible projection written above
and from the tensor-speed-safe scalar-tensor/DHOST restriction
\cite{Gleyzes2015,Langlois2016,Achour2016,Crisostomi2016}:
\[
G_4=\frac{M_{\rm Pl}^2}{2},\qquad G_5=0.
\]
The standard tensor-speed condition then gives
\[
\alpha_T=0,
\qquad
c_g=c.
\]
The spatial medium completion is static-silent for transverse-traceless modes
on the homogeneous background used here and adds no principal tensor correction
at this order, so
\[
\omega^2=c^2k^2.
\]
This is the local tensor-speed statement used in the paper; full cosmological
wave propagation, waveform damping, and catalog-level gravitational-wave tests
belong to the later propagation sector.

\section{Inertia and Mass}

This section fixes the energy normalization of the gravitational core. The same total localized dressed energy $E_0$ is read as the Noether inertial mass and as the far-zone refractive source charge. For a localized rest configuration, the standard relativistic normalization gives
\[
M_{\rm inert}=\frac{E_0}{c^2}.
\]
The weak gravitational source normalization gives the same mass:
\[
M_{\rm source}=\frac{E_0}{c^2}.
\]
Thus the leading identity used in this paper is
\[
M_{\rm inert}=M_{\rm source}=\frac{E_0}{c^2}.
\]
This is a minimally coupled rest-energy normalization, not an independent compact-body boost calculation.
The quantity $E_0$ is the total energy of the localized dressed state: the core and its exterior refractive dressing are counted once in the same normalization. Radiation and retarded dressing effects enter as higher-order response terms; the leading $F=Ma$ coefficient is fixed by the total rest energy. The full nonlinear compact-body Arnowitt--Deser--Misner (ADM)/Noether closure and finite-core dressing remain subsequent calculations.

\section{The Compact Phase Branch}

\subsection{Exponential Exterior}

The compact phase-spherical branch uses the reduced exterior static phase equation
\[
\left(r^2\phi'\right)'=0.
\]
With flat normalization at infinity,
\[
\phi=-\frac{r_s}{r}.
\]
The biconformal map is
\[
h=-\frac{\phi}{2},
\]
so that
\[
h=\frac{r_s}{2r}.
\]
The time and space conformal factors are
\[
e^{-2h},\qquad e^{2h}.
\]
The compact exterior metric is therefore
\[
ds^2=e^{-r_s/r}c^2dt^2
-e^{r_s/r}\left(dr^2+r^2d\Omega^2\right).
\]
It has $AB=1$, light index $n_{\rm light}=\sqrt{A/B}=e^{2h}$, and the weak limit $\gamma=1$. Here $B=e^{-2h}$ is the time factor and $A=e^{2h}$ is the spatial covariant factor.
Because $AB=1$, the curved static harmonic current
\[
\frac{\sqrt{-g}\,g^{rr}h'}{\sin\theta}
\]
reduces exactly to $-r^2h'$. Thus the curved phase equation on this branch is
the reduced radial equation written above.

The projected phase-pressure deficit source supporting this exterior is the
static reduction of the covariant operator written in the action section. On this branch
\[
H_\Delta=h,\qquad \omega_\Delta=1.
\]
The $H_\Delta$ equation has zero residual on this harmonic exterior.

\subsection{Projected Deficit Source}

The spatial medium is an internal label on the same branch. Use the
$SO(3)$-covariant map
\[
\phi^A=f(r)n^A(\theta,\varphi).
\]
Its radial and tangential eigenvalues are
\[
\lambda_r=\frac{f'^2}{A},\qquad
\lambda_t=\frac{f^2}{Ar^2},
\]
so
\[
I_3=\lambda_r\lambda_t^2
=e^{-6h}\frac{f'^2f^4}{r^4}.
\]
The regular identity exterior has $f=r$ and gives $I_3=e^{-6h}$. It serves as
the internal medium label. The active source law is the projected deficit
operator written below.

The radial action density per unit $\sin\theta$ is
\[
\mathcal L_{\Delta,r}
=\frac{\sqrt{-g}}{\sin\theta}
\left[
k\,\gamma^{rr}H_\Delta'^2
\right],
\qquad
k=\frac{\omega_\Delta}{8\pi G}.
\]
On the compact metric
\[
\frac{\sqrt{-g}}{\sin\theta}=r^2e^{2h},\qquad
\gamma^{rr}=e^{-2h}=\frac{1}{A}.
\]
The $H_\Delta$ Euler equation is therefore
\[
-2k\left(r^2H_\Delta'\right)'=0.
\]
For $H_\Delta=h=r_s/(2r)$ one has
\[
\left(r^2h'\right)'=0,
\]
so the deficit equation is solved. Since $f$ is absent from
$\mathcal L_{\Delta,r}$ as an active source variable, its source contribution is
$E_f=0$ on this exterior branch.
Therefore
\[
Z_\Delta^\perp=Z_\perp=\frac{(h')^2}{A}
=\frac{r_s^2e^{-r_s/r}}{4r^4},
\]
and
\[
\mathcal L_\Delta^\perp=\frac{Z_\Delta^\perp}{8\pi G}
=\frac{Z_\perp}{8\pi G}.
\]
Define
\[
\Delta_P=\frac{Z_\Delta^\perp}{8\pi G}
=\frac{Z_\perp}{8\pi G}
=\frac{r_s^2e^{-r_s/r}}{32\pi G r^4}.
\]
In the strong branch the deficit source is the same type of quadratic
pressure channel: the active source of the geometry contains the projected
square of the phase-index gradient $Z_\Delta^\perp$.
The active stress components are
\[
\Theta^t{}_t=-\Delta_P,\qquad
\Theta^r{}_r=\Delta_P,\qquad
\Theta^\theta{}_\theta=\Theta^\phi{}_\phi=-\Delta_P.
\]
With $p_r=-\Theta^r{}_r$ and $p_t=-\Theta^\theta{}_\theta$, this gives
\[
p_t-p_r=2\Delta_P.
\]
The Einstein tensor of the exponential exterior has the mixed profile
\[
G^t{}_t=-D,\qquad
G^r{}_r=D,\qquad
G^\theta{}_\theta=G^\phi{}_\phi=-D,
\]
with
\[
D=8\pi G\Delta_P.
\]
Equivalently,
\[
D=8\pi G\Delta_P=\frac{r_s^2e^{-r_s/r}}{4r^4}.
\]
Thus
\[
G^\mu{}_\nu-8\pi G\Theta^\mu{}_\nu=0.
\]
With this mixed-index convention, $D=8\pi G\Delta_P$ gives $G^t{}_t=-D=8\pi G\Theta^t{}_t$. The effective-fluid notation $\rho_a=-\Delta_P$ used below is the active deficit contribution of the medium, and the field-equation residual is read with this active contrast.
The compact-branch static consistency chain is then: $(r^2\phi')'=0$ fixes $\phi=-r_s/r$; the biconformal map gives $h=r_s/(2r)$ and the exponential metric; the independent deficit field has $H_\Delta=h$ and zero harmonic residual; $\mathcal L_\Delta^\perp=Z_\Delta^\perp/(8\pi G)$ gives the active source written above; with this source the diagonal field-equation residuals vanish and $p_t-p_r=2\Delta_P$. This is the static exterior equation of the compact branch. The finite-core selector is the load/source rule derived below.

\subsection{Raw-Invariant Diagnostic}

The tensor obstruction of the raw absolute-invariant $F_{\rm min}$ compact
insertion is explicit. Set
\[
w=e^{r_s/r}.
\]
If one uses the raw absolute invariants on the compact identity medium branch,
one has
\[
Y=w,\qquad \lambda_r=\lambda_t=\frac{1}{w}.
\]
On the physical Solar-family coefficient slice used for the weak branch, the
raw structural stress of $F_{\rm min}$ gives
\[
\frac{\Theta_F^t{}_t}{M_*^4c_{Y2}}
=\frac{(w-1)(3w^4-5w^3+w^2-23w+16)}{w^3},
\]
\[
\frac{\Theta_F^r{}_r}{M_*^4c_{Y2}}
=\frac{\Theta_F^\theta{}_\theta}{M_*^4c_{Y2}}
=\frac{\Theta_F^\phi{}_\phi}{M_*^4c_{Y2}}
=-\frac{(w-1)(w^4-7w^3-5w^2+3w+16)}{w^3}.
\]
For example,
\[
\left.\frac{\Theta_F^r{}_r}{M_*^4c_{Y2}}\right|_{w=2}=\frac{19}{4}\neq0.
\]
Thus adding this raw absolute-invariant stress to the balanced projected
deficit source would give the residual
\[
R_{\rm raw}^\mu{}_\nu
=G^\mu{}_\nu-8\pi G(\Theta_\Delta^\mu{}_\nu+\Theta_F^\mu{}_\nu)
=-8\pi G\,\Theta_F^\mu{}_\nu,
\]
which is nonzero. No constant retuning of the projected deficit load absorbs
this tensor, because the raw radial and tangential $F_{\rm min}$ stresses have
the same branch value while the compact geometry requires opposite
radial/tangential signs.

The compact action uses the phase-normalized strain variables instead:
\[
\widehat Y=e^{-2h}Y=1,\qquad
\widehat\lambda_r=e^{2h}\lambda_r=1,\qquad
\widehat\lambda_t=e^{2h}\lambda_t=1.
\]
Because the $F_{\rm min}$ unit state is tadpole-free, this gives
\[
\widehat F_{\rm min}=0,\qquad
\widehat\Theta_F^\mu{}_\nu=0,\qquad
E_f[\widehat F_{\rm min}]=0.
\]
The compact field equation is therefore evaluated with the active projected
deficit stress,
\[
G^\mu{}_\nu-8\pi G\Theta_\Delta^\mu{}_\nu=0,
\]
with the phase-normalized $F_{\rm min}$ action varied from the start. The
nonzero raw tensor above is the check of the absolute-invariant compact
insertion; the phase-normalized compact action has zero stress on the pure-phase
exterior.

Here the exponential form is derived from the compact phase-spherical branch:
the active projected deficit stress of the phase-spatial medium supplies the
exterior source.

The active exterior contribution alone has
\[
\rho_a=-\Delta_P,\qquad
p_{r,a}=-\Delta_P,\qquad
p_{t,a}=\Delta_P.
\]
This is the active contrast that sources the geometry. The two null-energy
sums of this source are
\[
\rho_a+p_{r,a}=-2\Delta_P,\qquad
\rho_a+p_{t,a}=0.
\]
Thus the deficit is read directly from the radial null sum:
\[
\Delta_P=-\frac{1}{2}\left(\rho_a+p_{r,a}\right).
\]
On the full analytic template the formal maximum of the deficit pressure is
\[
\Delta_P^{\rm max}=\frac{8e^{-4}}{\pi G r_s^2}.
\]
This point lies at $r=r_s/4$. The exterior traversable-wormhole domain begins
at the throat $r=r_s/2$. Therefore on the physical exterior $r\ge r_s/2$ the
deficit decreases outward, and the exterior maximum is the throat value
\[
\Delta_P^{\rm throat}=\frac{e^{-2}}{2\pi G r_s^2}.
\]
Since $\Delta_P>0$, the active exterior source violates the radial null energy condition (NEC), while
the transverse NEC is saturated. This violation is the algebraic signature of
the deficit: the standard effective-source language reads the phase-pressure
depletion of the base medium as a negative radial NEC sum.
This NEC check is applied to the active projected-deficit source
$\Theta_\Delta^\mu{}_\nu$. The compact $F_{\rm min}$ action is
phase-normalized and has zero compact-branch stress.

A homogeneous base background belongs to a different field-equation sector.
Adding it as a gravitating source changes the metric and requires a new
exterior derivation. The strong-branch result in this paper is therefore the
active-contrast exponential exterior in which the radial NEC sign is the direct
indicator of the phase-pressure deficit.

\subsection{Response, Energy Conditions, and Stiffness}

The same deficit admits a local mass-response balance. If a source of bare
amplitude $S>0$ is weakened by a deficit amplitude $q\ge0$ as
$S_{\rm eff}(q)=S e^{-\chi q}$, $\chi>0$, then the self-consistent balance
$q=S e^{-\chi q}$ gives
\[
q_\star=\frac{W(\chi S)}{\chi},\qquad W(x)e^{W(x)}=x.
\]
Linear stability follows from
\[
\left.\frac{d}{dq}\left(S e^{-\chi q}-q\right)\right|_{q=q_\star}
=-\chi q_\star-1<0.
\]
This fixed point supplies the local response rule; curved-background
perturbative stability is the separate principal-symbol and compact-spectrum
problem.

The static stiffness of the exterior source is positive:
\[
\delta^2\mathcal L_\Delta^\perp
=\frac{e^{-r_s/r}}{8\pi G}
\left[
(\delta h_r)^2
+\frac{(\delta h_\theta)^2}{r^2}
+\frac{(\delta h_\phi)^2}{r^2\sin^2\theta}
\right].
\]
This second-order operator in the independent deficit field $H_\Delta$ adds
positive static stiffness to the pressure-deficit channel. The propagating
hyperbolic modes are controlled by the full $F_{\rm min}$ medium sector.

\subsection{Static Geodesic Predictions}

The exponential exterior has no finite-radius event horizon in the static chart. Its strong-field use is therefore tied to tests of horizonless ultracompact objects. For such objects, a photon sphere has both a shadow-scale role and a dynamical-stability role: the light-ring structure is a stability criterion in its own right \cite{Cunha2017}, and the ringdown/echo channel gives a gravitational-wave discriminator from black holes \cite{Cardoso2016}. For horizonless light-ring objects, the quasinormal-mode (QNM)/echo channel is an open astrophysical-viability risk and part of the strong-field program itself. The result derived here is the static geometric core. Its astrophysical use proceeds through finite-core dynamics, the rotating exterior, QNM/echo analysis, and Event Horizon Telescope (EHT), next-generation Event Horizon Telescope (ngEHT), and Black Hole Explorer (BHEX) ray tracing.
For a matched exponential exterior, the numbers below depend only on the metric;
the finite-core selector fixes applicability, not the geodesic algebra.

Branch selection is set by the finite-core load and source amplitude. For an extended
weak object such as the Sun,
$C=r_s/R_{\rm source}\ll 1$; the diffuse medium stress and radial deformation
remain in the active-source exterior equation, and the weak-Solar branch gives
$q_{2{\rm PN}}=7/4$. For an ultradense compact object, the phase-pressure
deficit supplies the leading exterior source. The reduced exterior phase equation
$(r^2\phi')'=0$, the biconformal map, and the projected deficit source
$\mathcal L_\Delta^\perp$ define the exponential exterior. In this branch the
phase-normalized $F_{\rm min}$ sector has zero stress on the pure-phase exterior, and the
diagonal field-equation residuals vanish. A smooth $C^2$ exterior matching is
the condition under which no additional thin-shell stress is introduced at the
boundary.

The finite-core selector is algebraic at the static matching level.  Define the
load $Q=r_s/r_c$ and the input compactness $Q_0$.  Surface-clock matching
equates the volume-deficit clock to the compact surface lapse, and the same
radius-shrink load ledger gives
\[
e^{-q_v/3}=e^{-Q/2}\quad\Rightarrow\quad q_v=\frac{3Q}{2},\qquad
Q_0=Qe^{-Q/2}.
\]
The equation for $Q_0$ alone has two Lambert branches,
\[
Q_\pm=-2W_{0,-1}\!\left(-\frac{Q_0}{2}\right),
\]
so $Q_0$ by itself is not the branch selector.  The local deficit feedback is
\[
q_v=Se^{-\chi q_v},\qquad q_\star=\frac{W(\chi S)}{\chi},
\]
and hence
\[
Q_\star=\frac{2W(\chi S)}{3\chi}.
\]
Equivalently,
\[
S(Q)=\frac{3Q}{2}e^{3\chi Q/2},\qquad
\frac{dS}{dQ}=\frac34(3\chi Q+2)e^{3\chi Q/2}>0.
\]
Thus the finite-core source amplitude $S$ selects a unique load $Q$, not a
second root at fixed $Q_0$. The $Q_0$ fold point $Q=2$ maps to
\[
S_{\rm fold}=3e^{3\chi},
\]
but because $S(Q)$ is monotone this is a load marker in the $S$ variable, not a
bistable fold. For the finite-core horizonless exterior used by the geodesic
markers the required window is
\[
1<Q_\star<2,\qquad
\frac32 e^{3\chi/2}<S<3e^{3\chi},
\]
equivalently
\[
\frac{r_s}{2}<r_c<r_s.
\]
The throat is then inside the finite core while the photon sphere is exterior.
Within the static truncation this is the object class assigned to the loaded
compact $H=h$ exterior; object-specific values of $S$ and $\chi$ belong to the
finite-core equation of state or collapse history.

For null rays define
\[
B_{\rm ph}(r)=\frac{e^{-2r_s/r}}{r^2}.
\]
The exponential-metric geodesic calculation is standard
\cite{Boonserm2018}; in the present normalization the circular null condition
is checked in one line,
\[
B_{\rm ph}'(r)=\frac{2e^{-2r_s/r}(r_s-r)}{r^4}=0,
\]
and gives
\[
r_{\rm ph}=r_s.
\]
For the static spherical optical potential, the critical impact parameter is
$b_c=1/\sqrt{B_{\rm ph}(r_{\rm ph})}$; hence
\[
b_c=e\,r_s.
\]
For Schwarzschild at the same mass, $b_c^{\rm GR}=3\sqrt3\,r_s/2$.
Relative to Schwarzschild,
\[
\frac{b_c}{b_c^{\rm GR}}
=\frac{2e}{3\sqrt3}=1.0463.
\]
Thus the compact static branch gives the conditional shadow-scale marker
\[
\Delta b_c=+4.63\%.
\]

For massive circular orbits the effective potential is
\[
V_{\rm eff}=e^{-r_s/r}
+\frac{L^2e^{-2r_s/r}}{r^2}.
\]
Solving the standard circular-orbit and marginal-stability conditions for this
potential leaves the quadratic
\[
r^2-3r_s r+r_s^2=0,
\]
whose exterior root gives
\[
r_{\rm ISCO}
=\frac{3+\sqrt5}{2}\,r_s
=\varphi^2r_s.
\]
The orbital frequency is
\[
\Omega^2=
\frac{c^2 r_s e^{-2r_s/r}}{r^2(2r-r_s)}.
\]
At the ISCO, using the Schwarzschild normalization
$(\Omega_{\rm ISCO}^{\rm GR})^2=c^2/(54r_s^2)$ at the same mass,
\[
\frac{f_{\rm ISCO}}{f_{\rm ISCO}^{\rm GR}}
=\left[
\frac{54e^{-2/\varphi^2}}
{\varphi^4(2\varphi^2-1)}
\right]^{1/2}
=0.931.
\]
The conditional compact static branch markers are therefore the pair
\[
b_c=e\,r_s,\qquad r_{\rm ISCO}=\varphi^2r_s,
\]
with the ISCO frequency reduced by about $6.9\%$ relative to Schwarzschild at the same mass.

\begin{figure}[t]
\centering
\includegraphics[width=\textwidth]{figures/refg_geodesic_structure.pdf}
\caption{Strong-field geodesic diagnostics in units
$x=r/r_s$.  Panel (a) plots the null optical potential
$B_{\rm ph}=e^{-2r_s/r}/r^2$ for the RefG compact exterior and
$B_{\rm ph}^{\rm GR}=(1-r_s/r)/r^2$ for Schwarzschild, showing
$r_{\rm ph}=r_s$, $b_c=e\,r_s$ and
$b_c/b_c^{\rm GR}=2e/(3\sqrt3)=1.0463$.  Panel (b) plots the
massive effective potential
$V_{\rm eff}=e^{-r_s/r}+L^2e^{-2r_s/r}/r^2$ together with the
Schwarzschild comparison; the circular-orbit and marginal-stability conditions give
$r_{\rm ISCO}=\varphi^2r_s$ for RefG and $3r_s$ for Schwarzschild,
with $f_{\rm ISCO}/f_{\rm ISCO}^{\rm GR}=0.931$.}
\label{fig:geodesic-structure}
\end{figure}

\section{Falsification Criteria}

This RefG core is constrained by the following checks.

First, the local perturbation sector must have a healthy phase kinetic term and
a completed scalar-longitudinal principal symbol. The nondegenerate
bare-minimal window is
\[
c_{Y2}>0,\qquad -6c_{Y2}<\Sigma<-2c_{Y2},
\]
while the physical Solar condition sits on its longitudinal boundary
$K_{\pi\pi}=0$ and is accepted only after the dynamic completion. The final
determinant must have real, positive, non-superluminal scalar-longitudinal
roots,
\[
0<s_i\le 1.
\]
On the Solar $2$PN branch, the pre-completion $F_{\rm min}+C_6/Z$ determinant
derived in the local-stability section gives
$\det M_\odot(s)=5(s-1)^2$ at
$c_{Y2}=1$, $\lambda_6=1/4$, and $c_Z=1$. In this block $p_0=p_2$, so the
pre-completion determinant diagnoses the overconstraint.
The falsification filter is
therefore applied to the completed dynamic-channel criterion,
\[
B_{\rm new}=c_Z+\epsilon_B,\qquad
M_{\rm new}=8c_{Y2}-4\lambda_6-c_Z+\epsilon_M,
\]
with $A$, $C$, and $D$ kept at their static-core values.
$\lambda_6$, $c_Z$, $\epsilon_B$, and $\epsilon_M$ are EFT response
coefficients; the local-stability section gives the representative determinant
and roots.  On the same slice the open allowed region is
\[
26+\epsilon_B+10\sqrt{1+\epsilon_B}<(6+\epsilon_M)^2<36+6\epsilon_B,
\]
with width $5(\sqrt{1+\epsilon_B}-1)^2$.  The representative point
$\epsilon_B=1$, $\epsilon_M=\sqrt{83/2}-6$ has roots $0.5649$ and $0.8851$.
The falsification object is the open window itself; the displayed
representative point is one healthy completion, not a unique prediction for the
response coefficients.
Thus the scalar sector is part of the
falsification filter: a ghost, a complex root, a degenerate hyperbolic symbol,
or a superluminal scalar-longitudinal root rejects the parameter branch.

Second, the weak-Solar branch must keep
\[
\gamma=\beta=1,\qquad
\alpha_1=\alpha_2=\alpha_3=0,
\qquad
q_{2{\rm PN}}=\frac74.
\]

Third, tensor waves must remain luminal:
\[
c_g=c.
\]

Fourth, the compact selector must put the finite core in the compact exterior
window:
\[
Q_\star=\frac{2W(\chi S)}{3\chi},\qquad
1<Q_\star<2,
\]
equivalently
\[
\frac32 e^{3\chi/2}<S<3e^{3\chi},\qquad
\frac{r_s}{2}<r_c<r_s.
\]
If this condition fails, the finite-core object does not realize the compact
exterior window used by the geodesic markers. When it holds, the compact phase
branch gives conditional static branch markers:
\[
r_{\rm ph}=r_s,\qquad
b_c=e\,r_s,\qquad
r_{\rm ISCO}=\varphi^2r_s,\qquad
\frac{f_{\rm ISCO}}{f_{\rm ISCO}^{\rm GR}}=0.931.
\]
The shadow-scale shift is
\[
+4.63\%.
\]

The observational separation of this branch requires percent-level shadow-scale control, a rotating compact model, and plasma/accretion ray tracing for EHT, ngEHT, and BHEX-type tests.

\section{Conclusion}

RefG is an effective phase-spatial framework for gravity. Its source is the active stress of the medium, read through the source-index chain
\[
\Pi_\Delta'\rightarrow h_{\rm eff}\rightarrow n_{\rm eff}\rightarrow a_r,
\]
which gives Newton's law in the weak spherical limit. In the Solar branch the same action admits $\gamma=\beta=1$, the local preferred-frame source condition $K_{\rm pf}=0$, and the physical weak-Solar coefficient $q_{2{\rm PN}}=7/4$. In the tensor sector the principal wave speed has $c_g=c$.

The compact phase-spherical branch gives a static exponential exterior with no
finite-radius Killing horizon, supported by the projected phase-pressure
deficit stress and by the reduced exterior phase equation. Its $F_{\rm min}$
sector is evaluated on phase-normalized strain invariants; on the pure-phase
exterior this gives zero $F_{\rm min}$ metric stress and zero $\phi^A$ Euler
residual. The same active source obeys
\[
\Delta_P=-\frac{1}{2}\left(\rho_a+p_{r,a}\right),
\]
so the standard effective-source language directly reads the RefG
phase-pressure deficit. Its conditional static null and timelike geodesic
markers are
\[
b_c=e\,r_s,\qquad
r_{\rm ISCO}=\varphi^2r_s,
\]
with a $+4.63\%$ shadow-scale shift and
\[
f_{\rm ISCO}=0.931\,f_{\rm ISCO}^{\rm GR}.
\]
The finite-core source selector fixes the branch load by
\[
Q_\star=\frac{2W(\chi S)}{3\chi},
\]
with finite-core exterior applications requiring
$1<Q_\star<2$, equivalently $r_s/2<r_c<r_s$. Direct comparison with
EHT-class images belongs to the rotating exterior and ray-tracing problem.

The present paper gives the self-contained gravitational core. The next steps are object-specific finite-core equations of state, the rotating compact solution, quasinormal-mode and echo stability, and observational modeling of shadows, rings, and accretion environments. Those steps determine the direct comparison with EHT, ngEHT, and BHEX data.

\section*{Data availability statement}

No observational or experimental datasets were generated or analysed in this
study. The results are based on analytic derivations presented in the article
and on cited standard results.

\begin{thebibliography}{99}

\bibitem{Horndeski1974}
G. W. Horndeski, Second-order scalar-tensor field equations in a four-dimensional space, \textit{International Journal of Theoretical Physics} \textbf{10}, 363--384 (1974).

\bibitem{Gubitosi2013}
G. Gubitosi, F. Piazza, and F. Vernizzi, The effective field theory of dark energy, \textit{Journal of Cosmology and Astroparticle Physics} \textbf{2013}, 032 (2013).

\bibitem{BelliniSawicki2014}
E. Bellini and I. Sawicki, Maximal freedom at minimum cost: linear large-scale structure in general modifications of gravity, \textit{Journal of Cosmology and Astroparticle Physics} \textbf{2014}, 050 (2014).

\bibitem{Abbott2017GW}
B. P. Abbott et al., GW170817: Observation of gravitational waves from a binary neutron star inspiral, \textit{Physical Review Letters} \textbf{119}, 161101 (2017).

\bibitem{Abbott2017GRB}
B. P. Abbott et al., Gravitational waves and gamma-rays from a binary neutron star merger: GW170817 and GRB 170817A, \textit{Astrophysical Journal Letters} \textbf{848}, L13 (2017).

\bibitem{Endlich2013}
S. Endlich, A. Nicolis, and J. Wang, Solid inflation, \textit{Journal of Cosmology and Astroparticle Physics} \textbf{2013}, 011 (2013).

\bibitem{Nicolis2015}
A. Nicolis, R. Penco, F. Piazza, and R. A. Rosen, More on gapped Goldstones at finite density: more gapped Goldstones, \textit{Journal of High Energy Physics} \textbf{2015}, 055 (2015).

\bibitem{Will1993}
C. M. Will, \textit{Theory and Experiment in Gravitational Physics}, Cambridge University Press (1993).

\bibitem{Will2014}
C. M. Will, The confrontation between general relativity and experiment, \textit{Living Reviews in Relativity} \textbf{17}, 4 (2014).

\bibitem{Bertotti2003}
B. Bertotti, L. Iess, and P. Tortora, A test of general relativity using radio links with the Cassini spacecraft, \textit{Nature} \textbf{425}, 374--376 (2003).

\bibitem{Creminelli2017}
P. Creminelli and F. Vernizzi, Dark energy after GW170817 and GRB170817A, \textit{Physical Review Letters} \textbf{119}, 251302 (2017).

\bibitem{Sakstein2017}
J. Sakstein and B. Jain, Implications of the neutron star merger GW170817 for cosmological scalar-tensor theories, \textit{Physical Review Letters} \textbf{119}, 251303 (2017).

\bibitem{Deffayet2009}
C. Deffayet, G. Esposito-Farese, and A. Vikman, Covariant Galileon, \textit{Physical Review D} \textbf{79}, 084003 (2009).

\bibitem{Kobayashi2011}
T. Kobayashi, M. Yamaguchi, and J. Yokoyama, Generalized G-inflation: inflation with the most general second-order field equations, \textit{Progress of Theoretical Physics} \textbf{126}, 511--529 (2011).

\bibitem{Gleyzes2015}
J. Gleyzes, D. Langlois, F. Piazza, and F. Vernizzi, Healthy theories beyond Horndeski, \textit{Physical Review Letters} \textbf{114}, 211101 (2015).

\bibitem{Langlois2016}
D. Langlois and K. Noui, Degenerate higher derivative theories beyond Horndeski: evading the Ostrogradski instability, \textit{Journal of Cosmology and Astroparticle Physics} \textbf{2016}, 034 (2016).

\bibitem{Achour2016}
J. Ben Achour, D. Langlois, and K. Noui, Degenerate higher order scalar-tensor theories beyond Horndeski and disformal transformations, \textit{Physical Review D} \textbf{93}, 124005 (2016).

\bibitem{Crisostomi2016}
M. Crisostomi, M. Hull, K. Koyama, and G. Tasinato, Horndeski: beyond, or not beyond?, \textit{Journal of Cosmology and Astroparticle Physics} \textbf{2016}, 038 (2016).

\bibitem{Cheung2008}
C. Cheung, P. Creminelli, A. L. Fitzpatrick, J. Kaplan, and L. Senatore, The effective field theory of inflation, \textit{Journal of High Energy Physics} \textbf{2008}, 014 (2008).

\bibitem{Endlich2014}
S. Endlich, B. Horn, A. Nicolis, and J. Wang, The squeezed limit of the solid inflation three-point function, \textit{Physical Review D} \textbf{90}, 063506 (2014).

\bibitem{Bordin2017}
L. Bordin, P. Creminelli, M. Mirbabayi, and J. Nore\~na, Solid consistency, \textit{Journal of Cosmology and Astroparticle Physics} \textbf{2017}, 004 (2017).

\bibitem{Kostelecky2022}
V. A. Kosteleck\'y and N. Russell, Data tables for Lorentz and CPT violation, \textit{Reviews of Modern Physics} \textbf{83}, 11 (2011); updated editions available at arXiv:0801.0287.

\bibitem{WillNordtvedt1972}
C. M. Will and K. Nordtvedt, Conservation laws and preferred frames in relativistic gravity. I. Preferred-frame theories and an extended PPN formalism, \textit{Astrophysical Journal} \textbf{177}, 757--774 (1972).

\bibitem{Damour1992}
T. Damour and G. Esposito-Farese, Tensor-multi-scalar theories of gravitation, \textit{Classical and Quantum Gravity} \textbf{9}, 2093--2176 (1992).

\bibitem{Damour1993}
T. Damour and G. Esposito-Farese, Nonperturbative strong-field effects in tensor-scalar theories of gravitation, \textit{Physical Review Letters} \textbf{70}, 2220--2223 (1993).

\bibitem{Damour1996}
T. Damour and G. Esposito-Farese, Tensor-scalar gravity and binary-pulsar experiments, \textit{Physical Review D} \textbf{54}, 1474--1491 (1996).

\bibitem{Boonserm2018}
P. Boonserm, T. Ngampitipan, A. Simpson, and M. Visser, Exponential metric represents a traversable wormhole, \textit{Physical Review D} \textbf{98}, 084048 (2018).

\bibitem{Cunha2017}
P. V. P. Cunha, E. Berti, and C. A. R. Herdeiro, Light-ring stability for ultracompact objects, \textit{Physical Review Letters} \textbf{119}, 251102 (2017).

\bibitem{Cardoso2016}
V. Cardoso, S. Hopper, C. F. B. Macedo, C. Palenzuela, and P. Pani, Gravitational-wave signatures of exotic compact objects and of quantum corrections at the horizon scale, \textit{Physical Review D} \textbf{94}, 084031 (2016).

\end{thebibliography}

\end{document}
